
\documentclass[a4paper,10pt]{report}

%\usepackage[landscape]{geometry}
\usepackage[cm-default]{fontspec}
\setromanfont{FreeSerif}
\setsansfont{FreeSans}
\setmonofont{FreeMono}
\usepackage[utf8x]{inputenc}
\usepackage{fontspec}
\usepackage{xunicode}
\usepackage{xltxtra}
\usepackage{xgreek}
\usepackage{amsmath}
\usepackage{unicode-math}
\usepackage{ulem}
\usepackage{color}
\usepackage{verbatim}
\usepackage{nopageno}
\usepackage{graphicx}
\usepackage{textpos}
\setlength{\TPHorizModule}{1cm}
\setlength{\TPVertModule}{1cm}
%\usepackage[colorgrid,texcoord]{eso-pic}
\usepackage[outline]{contour}
\usepackage{wrapfig}
\usepackage{url}
\usepackage{color}
\usepackage{tikz}
\usepackage{fancybox,fancyhdr}
\usepackage{subfigure}
\usepackage{pstricks}
\usepackage{epsfig}
\usepackage{multicol}
\usepackage{listings}
\usepackage{enumerate}
\usepackage{hyperref}
\hypersetup{
  bookmarks=true,
  bookmarksopen=true,
  pdfborder=false,
  pdfpagemode=UseNone,
  raiselinks=true,
  pdfhighlight={/P},
  colorlinks,
  citecolor=black,
  filecolor=black,
  linkcolor=black,
  urlcolor=black
}


\usepackage{fontspec}
\usepackage{xunicode}
\usepackage{xltxtra}


% Margins
%\setlength{\textwidth}{24cm}

%\setlength{\voffset}{0in}
%\setlength{\textheight}{6.8in}
% Colors
\definecolor{cbrown}{rgb}{0.49,0.24,0.07}
\definecolor{cmpez}{rgb}{0.92,0.88,0.79}
\definecolor{cmpez2}{rgb}{0.62,0.33,0.34}
\definecolor{cwhite}{rgb}{1,1,1}
\definecolor{cblack}{rgb}{0,0,0}
\definecolor{cred}{rgb}{0.9,0.15,0.15}
\definecolor{clblue}{rgb}{0.57,0.85,0.97}
\definecolor{corange}{rgb}{0.86,0.49,0.18}
\definecolor{clorange}{rgb}{0.95,0.84,0.61}
\definecolor{cyellow}{rgb}{1,0.95,0.16}
\definecolor{cgreen}{rgb}{0.69,0.8,0.31}
\definecolor{clbrown}{rgb}{0.78,0.67,0.43}
\definecolor{cmagenta}{rgb}{0.79,0.05,0.54}
\definecolor{cgray}{rgb}{0.76,0.73,0.63}
\definecolor{bluesite}{rgb}{0.05,0.43,0.67}
\definecolor{redsite}{rgb}{0.90,0.15,0.15}
\definecolor{backsite}{rgb}{0.07,0.07,0.07}
% ----------- from Costas -----------------------------

%---------------   Main Fonts    -----------------%   

% \setmainfont[Mapping=tex-text]{Calibri}
% \setmainfont[Mapping=tex-text]{Times New Roman}
%\setmainfont[Mapping=tex-text]{Droid Serif}
%\setmainfont[Mapping=tex-text]{Cambria}
%\setmainfont[Mapping=tex-text]{cm-unicode}
% \setmainfont[Mapping=tex-text]{Gentium}
% \setmainfont[Mapping=tex-text]{GFS Didot}
% \setmainfont[Mapping=tex-text]{Comic Sans MS}
 \setmainfont[Mapping=tex-text]{Ubuntu}
% \setmainfont[Mapping=tex-text]{Myriad Pro}
% \setmainfont[Mapping=tex-text]{CMU Concrete}
%\setmainfont[Mapping=tex-text]{DejaVu Sans}
%\setmainfont[Mapping=tex-text]{KerkisSans}
%\setmainfont[Mapping=tex-text]{KerkisCaligraphic}
% \setmainfont[Mapping=tex-text]{Segoe Print}
% \setmainfont[Mapping=tex-text]{Gabriola}

%------------------------------------------------%

% ------------ Mathematics Fonts ----------------%
\setmathfont{Asana-Math.ttf}
% -----------------------------------------------%
% Misc
\author{Κασωτάκη Ε. - Σμαραγδάκης Κ.}
\title{www.math24.gr}
% ---------------- Formation --------------------%


\setlength\topmargin{-0.7cm}
\addtolength\topmargin{-\headheight}
\addtolength\topmargin{-\headsep}
\setlength\textheight{26cm}
\setlength\oddsidemargin{-0.54cm}
\setlength\evensidemargin{-0.54cm}
%\setlength\marginparwidth{1.5in}
\setlength\textwidth{17cm}

\RequirePackage[avantgarde]{quotchap}
\renewcommand\chapterheadstartvskip{\vspace*{0\baselineskip}}
\RequirePackage[calcwidth]{titlesec}
\titleformat{\section}[hang]{\bfseries}
{\Large\thesection}{12pt}{\Large}[{\titlerule[0.9pt]}]
%--------------------------------------------------%
\usepackage{draftwatermark}
\SetWatermarkText{www.math24.gr}
\SetWatermarkLightness{0.9}
\SetWatermarkScale{0.6}
%--------------------------------------------------%


\usepackage{xcolor}
\usepackage{amsthm}
\usepackage{framed}
\usepackage{parskip}

\colorlet{shadecolor}{bluesite!20}

\newtheorem*{orismos}{Ορισμός}

\newenvironment{notation}
  {\begin{shaded}\begin{theorem}}
  {\end{theorem}\end{shaded}}




% Document begins
\begin{document}
\pagestyle{fancy}
\fancyhead{}
\fancyfoot{}
\renewcommand{\headrulewidth}{0pt}
\renewcommand{\footrulewidth}{0pt}


\fancyhead[LO,LE]{
 %\textblockcolor{backsite}
 %\begin{textblock}{5}(-2,-0.55)
  %\rule{0cm}{1cm}
 %\end{textblock}
 \textblockcolor{bluesite}
 \begin{textblock}{5}(-1.5,-0.55)
  \rule{0cm}{1cm}
 \end{textblock}
 %\textblockcolor{bluesite}
 %\begin{textblock}{14}(5.5,-0.55)
 % \rule{0cm}{1cm}
 %\end{textblock}
 \begin{textblock}{0}(-1,-0.25)
 \color{cwhite} \begin{Large}www.math24.gr\end{Large}
 \end{textblock}
\begin{textblock}{0}(16,-0.4)
 \color{cwhite} \includegraphics[height=1.5cm]{math24_logo.png}
 \end{textblock}
\textblockcolor{white}
\begin{textblock}{17}(-1,28)
 \color{backsite} \begin{small}Copyright \textcopyright 2011, Κασωτάκη Ε.(ikasotaki@gmail.com) - Σμαραγδάκης Κ.(kesmarag@gmail.com)\end{small}
 \end{textblock}
}
 \begin{shaded}
\begin{center}
\huge \textbf{Μαθηματικά Α' Γυμνασίου}\\
\end{center} 
\begin{center}
\textbf{Επαναληπτικές Ασκήσεις για το Κεφάλαιο 1 "Οι φυσικοί αριθμοί"} \hfill \textbf{}
\end{center}
%\subsection*{}
\end{shaded}
%\vspace{2em
%\begin{shaded}
%\begin{center}
%\huge \textbf{Επαναληπτικές Ασκήσεις για το Κεφάλαιο 1}\\
%Πρόσθεση ρητών αριθμών
%\end{center} 
%\textbf{Μαθηματικά Α' Γυμνασίου} \hfill \textbf{Ημερομηνία Παράδοσης : \hspace{2em} }
%\subsection*{Ονοματεπώνυμο :\hfill  \hspace{5em}}
%\end{shaded}
\vspace{2em}
\begin{itemize}
\item Πρόσθεση, αφαίρεση και πολλαπλασιασμός φυσικών αριθμών
\item Δυνάμεις φυσικών αριθμών
\item Ευκλείδεια διαίρεση - Διαιρετότητα
\item Χαρακτήρες διαιρετότητας - Μ.Κ.Δ - Ε.Κ.Π - Ανάλυση αριθμού σε γινόμενο πρώτων παραγόντων
\end{itemize}


\section*{Άσκηση 1  \hfill \small{}}
Να συμπληρώσετε τα παρακάτω κενά:
\begin{enumerate}[1)]
 \item Η ιδιότητα $α+β=β+α$ λέγεται ............................
 \item Το 1 όταν ................................ με έναν φυσικό αριθμό δεν τον μεταβάλλει.
 \item Για να πολλαπλασιάσουμε έναν φυσικό αριθμό επί 1000 γράφουμε στο τέλος του αριθμού ..........................
       μηδενικά.
 \item Οι δυνάμεις του 1, δηλαδή $1^{ν}$, είναι όλες ίσες με ....................
 \item Στους φυσικούς αριθμούς η τέλεια διαίρεση είναι πράξη αντίστροφη του ...................................
 \item Ο διαιρέτης μιας διαίρεσης δεν μπορεί να είναι ..............................
 \item Σε μια διαίρεση όταν ο διαιρετέος είναι 0 τότε το πηλίκο είναι (ισούται με) ..............................
 \item Κάθε φυσικός αριθμός .............................. τα πολλαπλάσιά του.
 \item Ένας αριθμός που έχει διαιρέτες μόνο τον εαυτό του και το 1 λέγεται .............................
 \item Ένας φυσικός αριθμός ........................... με το 5 αν λήγει σε 0 ή 5.
\end{enumerate}

\section*{Άσκηση 2  \hfill \small{}}
Να χαρακτηρίσετε καθεμία από τις παρακάτω προτάσεις με τη λέξη "\textbf{Σωστό}``, αν είναι σωστή ή  
"\textbf{Λάθος}`` , αν είναι λανθασμένη:
\begin{enumerate}[1)]
 \item Η διαφορά δύο περιττών αριθμών είναι πάντα περιττός αριθμός
 \item Η πράξη $50-(13-3)$ δίνει το ίδιο αποτέλεσμα με την πράξη $50-13-3$
 \item Για να πολλαπλασιάσουμε ένα φυσικό αριθμό με το 10000 γράφουμε στο τέλος του αριθμού 4 μηδενικά
 \item Η πράξη $17\cdot(3+10)$ δίνει το ίδιο αποτέλεσμα με την πράξη $17\cdot3+17\cdot10$
 \item Το $2^{4}$ ισούται με $8$
 \item $2+2+2+2=4\cdot2$
 \item $2^{3}=3^{2}$
 \item $4\cdot10^{2}+3\cdot10^{1}+4$ είναι το ανάπτυγμα του αριθμού $434$ σε δυνάμεις του $10$
 \item Ο διαιρέτης μιας διαίρεση δεν μπορεί να είναι 0
 \item Ο διαιρετέος μιας διαίρεση δεν μπορεί να είναι 0
 \item Η σχέση $22=4\cdot5+2$ είναι μια ευκλείδεια διαίρεση
 \item Ο αριθμός $3127$ διαιρείται με το $5$
 \item Ο αριθμός $9α+3$ διαιρείται με το $3$
 \item Ο αριθμός $20$ αναλύεται σε γινόμενο πρώτων παραγόντων ως $2^{2}\cdot 5$
 \item Ο αριθμός $200$ αναλύεται σε γινόμενο πρώτων παραγόντων ως $2\cdot 10^{2}$
 \item Το ΕΚΠ των 3 και 6 είναι το 18
 \item Ο ΜΚΔ των 12 και 24 είναι το 6
 \item Το ΕΚΠ των $2^{2}\cdot 3^{4}$ και $2^{3}\cdot 3^{3}$ είναι $2^{3}\cdot 3^{4}$
 \item O MKΔ των $2^{3}\cdot 3^{5}$ και $2^{2}\cdot 3^{7}$ είναι $2^{2}\cdot 3^{5}$
 \item Οι αριθμοί 270 και 135 έχουν μέγιστο κοινό διαιρέτη τον αριθμό 5
\end{enumerate}


\section*{Άσκηση 3  \hfill \small{}}
Να αντιστοιχίσετε κάθε στοιχείο της αριστερής στήλης με ένα στοιχείο της δεξιάς στήλης
\begin{itemize}
\begin{multicols}{2}
 \item ΕΚΠ(2,6)
 \item ΕΚΠ(8,32)
 \item ΕΚΠ(7,14)
 \item ΜΚΔ(100,1000)
 \item ΜΚΔ(5,32)
 \item ΜΚΔ(9,24)
 \item 32
 \item 100
 \item 3
 \item 6
 \item 14
 \item 1
\end{multicols}
\end{itemize}


\section*{Άσκηση 4  \hfill \small{}}
Να βρείτε το αποτέλεσμα για καθεμία από τις παρακάτω αριθμητικές παραστάσεις:
\begin{enumerate}[1)]
 \item $3\cdot5\cdot2^{2}+5^{2}\cdot3+10$
 \item $3\cdot(5\cdot2^{2}+5^{2}\cdot3)+10$
 \item $3\cdot(5\cdot2^{2}+5^{2}\cdot3+10)$
 \item $3\cdot5\cdot(2^{2}+5^{2}\cdot3)+10$
 \item $3\cdot5\cdot2^{2}+5^{2}\cdot(3+10)$
\end{enumerate}


\section*{Άσκηση 5  \hfill \small{}}
Να γράψετε τους παρακάτω αριθμούς σε αναπτυγμένη μορφή με χρήση των δυνάμεων του 10
\begin{enumerate}[1)]
 \item 420
 \item 567
\end{enumerate}


\section*{Άσκηση 6  \hfill \small{}}
Να βρείτε αν οι παρακάτω αριθμοί διαιρούνται με $2,5,9$
\begin{enumerate}[1)]
 \item 3301
 \item 4075
 \item 8010
 \item 1219
 \item 90
\end{enumerate}

\section*{Άσκηση 7  \hfill \small{}}
Να βρείτε το ΕΚΠ των παρακάτω αριθμών: 
\begin{enumerate}[1)]
 \item 4,10
 \item 3,5
 \item 6,13
 \item 12,16
 \item 13,26
\end{enumerate}

\section*{Άσκηση 8  \hfill \small{}}
Να βρείτε το ΜΚΔ των παρακάτω αριθμών: 
\begin{enumerate}[1)]
 \item 12,36
 \item 17,33
 \item 45,75
 \item 24,28
 \item 7,21
\end{enumerate}

\section*{Άσκηση 9  \hfill \small{}}
Να αναλυθούν οι αριθμοί 320 και 1320 σε γινόμενα πρώτων παραγόντων. Με την βοήθεια της ανάλυσης αυτής να βρεθούν 
ο ΜΚΔ και το ΕΚΠ αυτών των αριθμών

\section*{Άσκηση 10  \hfill \small{}}
Να αναλυθούν οι αριθμοί 256 και 960 σε γινόμενα πρώτων παραγόντων. Με την βοήθεια της ανάλυσης αυτής να βρεθούν 
ο ΜΚΔ και το ΕΚΠ αυτών των αριθμών

\end{document}
